\section{Delimitación y fundamentación del problema de investigación}
\label{sec:problema}

\subsection{Descripción técnica del problema}

La problemática central se define como la incapacidad de las actuales
plataformas de automatización y adquisición de datos industriales para
representar de forma consistente y precisa el estado interno de carga de la
tolva de recepción en un chancador primario minero.

El sistema físico es inherentemente estocástico y parcialmente observable.
Variables operacionales críticas —el volumen neto de mineral retenido en la
cavidad interna y, sobre todo, la tasa instantánea de extracción del
chancador— no pueden sensorizarse de manera directa debido a la hostilidad
ambiental y a las limitaciones geométricas del equipo. Los datos crudos
provistos por el SCADA ingresan corrompidos por el ruido del entorno, lo que
provoca que los modelos deterministas tradicionales, basados en umbrales fijos y
alarmas de nivel, fallen al capturar las perturbaciones dinámicas del proceso.

La delimitación del problema es deliberadamente estrecha. No se estudia la
totalidad del proceso de conminución ni el diseño del gemelo digital completo,
sino un componente específico y acotado: la capa lógica que media entre la
telemetría degradada y cualquier consumidor analítico posterior.

\subsection{Causas, consecuencias e impacto informático}

Las causas del problema corresponden a las limitaciones de precisión de los
instrumentos físicos industriales operando bajo vibración constante severa y
atmósfera con material particulado en suspensión. Se trata de una restricción
física de la medición, no de un defecto de diseño del sistema de control.

Las consecuencias a nivel de planta se traducen en detenciones no programadas
por atoro de mineral en la tolva, ineficiencia en el consumo de energía
eléctrica del motor principal y desgaste prematuro de los componentes mecánicos
del activo.

Desde la perspectiva de la ingeniería informática, sin embargo, el impacto es de
otra naturaleza y es el que interesa a este estudio. Las herramientas avanzadas
de optimización digital, los tableros analíticos y los algoritmos predictivos
característicos de la Industria 4.0 dependen por completo de la consistencia y
veracidad de sus entradas. Si el \emph{pipeline} de datos procesa entradas
corruptas, las decisiones automatizadas divergen y generan recomendaciones
subóptimas. El problema, formulado en términos informáticos, es que se está
construyendo una capa de inteligencia artificial industrial sobre un sustrato de
datos cuya fidelidad nadie ha garantizado.

Esta observación no es especulativa. La revisión sistemática de
\textcite{ebad2026} sobre sesenta y ocho estudios de gemelos digitales
impulsados por inteligencia artificial en minería identifica la integración de
datos y el procesamiento en tiempo real entre los desafíos centrales no
resueltos del área, junto con la escasez de validación industrial efectiva.

\subsection{Justificación y pertinencia del estudio}

Esta investigación se justifica por la necesidad de proveer una solución que
resuelva el vacío informacional de forma eficiente y de bajo costo.

A diferencia de los enfoques convencionales, que requieren inversiones
significativas en hardware de instrumentación o la reestructuración completa de
las redes SCADA existentes, este proyecto propone el diseño de un módulo de
software analítico autocontenido y desacoplado. Mediante la aplicación de
filtros estocásticos recursivos fundamentados en leyes físicas de conservación,
el módulo depura las señales en tiempo de operación. Esto permite resguardar la
consistencia de los datos directamente en la capa de software, habilitando la
viabilidad de los modelos analíticos de nivel superior de manera sostenible.

La pertinencia disciplinar del estudio se sostiene en tres planos. En el plano
\textbf{teórico}, contribuye a precisar la distinción conceptual entre un
\emph{digital shadow}, que se limita a replicar datos crudos, y un verdadero
\emph{digital twin}, que procesa y depura la incertidumbre subyacente
\parencite{sepasgozar2021}. En el plano \textbf{tecnológico}, entrega un
componente de software reproducible, verificable y reutilizable sobre distintos
subsistemas del mismo activo. Y en el plano \textbf{contextual}, aborda un
proceso que constituye el eslabón inicial de la cadena productiva de la
principal industria exportadora de Chile, en un momento en que la adopción de
gemelos digitales en la industria de minerales se encuentra en expansión
documentada \parencite{qu2023}.
